\chapter{第一段可维护程序}

\section{问题入口：统计一组成绩}

第一段程序不要只写 \code{Hello, world}。我们从一个小需求开始：读入若干成绩，输出人数、总分、平均分、最高分和最低分。这个需求足够小，但已经覆盖输入、容器、循环、函数和错误处理。

最直接的写法可能是把所有逻辑塞进 \code{main}：

\begin{lstlisting}[language=C++,caption={所有逻辑挤在 main 里}]
#include <iostream>
#include <vector>

int main()
{
    std::vector<int> scores;
    int value = 0;
    while (std::cin >> value) {
        scores.push_back(value);
    }

    int sum = 0;
    int min_score = scores[0];
    int max_score = scores[0];
    for (int score : scores) {
        sum += score;
        if (score < min_score) {
            min_score = score;
        }
        if (score > max_score) {
            max_score = score;
        }
    }

    std::cout << scores.size() << '\n';
    std::cout << sum / scores.size() << '\n';
    std::cout << min_score << " " << max_score << '\n';
}
\end{lstlisting}

它能在普通输入下运行，但有三个问题。第一，空输入时 \code{scores[0]} 越界。第二，平均分用了整数除法，\code{95} 和 \code{96} 的平均值会被截断。第三，读取、统计和输出混在一起，后面不好测试。

\section{把数据结果变成类型}

先设计一个结果类型。统计结果不是五个散落的变量，而是一个概念：\code{ScoreSummary}。

\begin{lstlisting}[language=C++,caption={用结构体表达统计结果}]
#include <cstddef>

struct ScoreSummary {
    std::size_t count = 0;
    int sum = 0;
    int min_score = 0;
    int max_score = 0;
    double average = 0.0;
};
\end{lstlisting}

这个结构体没有行为，只是承载结果。它的默认值也有意义：当没有成绩时，\code{count} 是 \code{0}，平均值是 \code{0.0}。但是否允许空输入，不应该靠默认值偷偷决定，而应该由统计函数明确处理。

\section{函数边界：只观察，不拥有}

统计函数只需要观察一段成绩，不需要复制它。用 \code{std::span<const int>} 可以表达“我看一段连续整数，但不拥有它”。

\begin{lstlisting}[language=C++,caption={统计函数使用 span 观察数据}]
#include <algorithm>
#include <numeric>
#include <optional>
#include <span>
#include <vector>

std::optional<ScoreSummary> summarize_scores(std::span<const int> scores)
{
    if (scores.empty()) {
        return std::nullopt;
    }

    const auto [min_it, max_it] = std::minmax_element(scores.begin(), scores.end());
    const int sum = std::accumulate(scores.begin(), scores.end(), 0);

    ScoreSummary summary;
    summary.count = scores.size();
    summary.sum = sum;
    summary.min_score = *min_it;
    summary.max_score = *max_it;
    summary.average = static_cast<double>(sum) / static_cast<double>(scores.size());
    return summary;
}
\end{lstlisting}

\code{std::optional} 表达“可能没有结果”。这比返回一个看似正常但其实无意义的 \code{ScoreSummary} 更诚实。调用者必须处理空输入。

\section{输入和输出留在边缘}

现在 \code{main} 只做三件事：读入、调用核心逻辑、输出结果。

\begin{lstlisting}[language=C++,caption={把核心逻辑从 main 中拆出来}]
#include <iostream>

std::vector<int> read_scores()
{
    std::vector<int> scores;
    int value = 0;
    while (std::cin >> value) {
        scores.push_back(value);
    }
    return scores;
}

int main()
{
    const std::vector<int> scores = read_scores();
    const std::optional<ScoreSummary> summary = summarize_scores(scores);
    if (!summary.has_value()) {
        std::cerr << "no scores\n";
        return 1;
    }

    std::cout << "count: " << summary->count << '\n';
    std::cout << "sum: " << summary->sum << '\n';
    std::cout << "average: " << summary->average << '\n';
    std::cout << "range: " << summary->min_score
              << "..." << summary->max_score << '\n';
}
\end{lstlisting}

这段代码的重点不是更长，而是责任更清楚。读输入的函数可以单独替换，统计函数可以用固定数组测试，输出格式可以改而不碰统计逻辑。

\begin{keyidea}
好程序不是从“少写几行”开始，而是从“每段代码只负责一件事”开始。核心逻辑尽量不直接依赖终端、文件和全局状态。
\end{keyidea}

\section{本章边界检查}

这个小程序至少要检查四类输入：

\begin{itemize}
  \item 空输入：应返回 \code{std::nullopt}。
  \item 单个成绩：最小值和最大值相同。
  \item 多个成绩：平均分不能被整数除法截断。
  \item 非法文本：当前读入逻辑会在失败处停止，这是否符合需求要由调用场景决定。
\end{itemize}

从第一段程序开始，就把“能跑”和“可解释、可测试、可维护”分开。后者才是工程里的 \cpp。

\section{一步一步推导统计函数}

为什么统计函数要先处理空输入？从最小反例看最清楚。输入为空时，根本不存在最小值、最大值和平均值。如果代码为了省事把它们设成 \code{0}，调用者就会误以为真的统计到了一个值为 \code{0} 的成绩。于是函数的返回类型不能只是 \code{ScoreSummary}，还要能表达“没有结果”。

再看单元素输入：

\begin{lstlisting}[language=C++,caption={单元素时每个字段都能推导出来}]
const std::vector<int> scores{91};
// count = 1
// sum = 91
// min_score = 91
// max_score = 91
// average = 91.0
\end{lstlisting}

单元素不是特殊麻烦，而是最好的基准。只要它成立，多元素就可以理解成“把同一套规则重复应用”。扫描第一个元素时，最小值和最大值都等于它；扫描后续元素时，只在发现更小或更大时更新边界。标准库的 \code{std::minmax_element} 把这个扫描过程封装好了，但你仍然要知道它在做什么。

\section{完整文件版本}

前面的代码片段可以合成一个可运行文件。入门阶段不要只看片段，至少要能把它们合起来编译：

\begin{lstlisting}[language=C++,caption={完整成绩统计程序}]
#include <algorithm>
#include <iostream>
#include <numeric>
#include <optional>
#include <span>
#include <vector>

struct ScoreSummary {
    std::size_t count = 0;
    int sum = 0;
    int min_score = 0;
    int max_score = 0;
    double average = 0.0;
};

std::optional<ScoreSummary> summarize_scores(std::span<const int> scores)
{
    if (scores.empty()) {
        return std::nullopt;
    }

    const auto [min_it, max_it] = std::minmax_element(scores.begin(), scores.end());
    const int sum = std::accumulate(scores.begin(), scores.end(), 0);

    return ScoreSummary{
        .count = scores.size(),
        .sum = sum,
        .min_score = *min_it,
        .max_score = *max_it,
        .average = static_cast<double>(sum) / static_cast<double>(scores.size()),
    };
}

std::vector<int> read_scores()
{
    std::vector<int> scores;
    int value = 0;
    while (std::cin >> value) {
        scores.push_back(value);
    }
    return scores;
}

int main()
{
    const std::vector<int> scores = read_scores();
    const std::optional<ScoreSummary> summary = summarize_scores(scores);
    if (!summary.has_value()) {
        std::cerr << "no scores\n";
        return 1;
    }

    std::cout << "count: " << summary->count << '\n';
    std::cout << "sum: " << summary->sum << '\n';
    std::cout << "average: " << summary->average << '\n';
    std::cout << "range: " << summary->min_score
              << "..." << summary->max_score << '\n';
}
\end{lstlisting}

运行时可以这样喂输入：

\begin{lstlisting}[language=bash,caption={用管道提供输入}]
echo "90 80 100" | ./score
\end{lstlisting}

如果没有任何输入，程序返回错误码 \code{1}。如果有输入，核心统计逻辑完全不关心输入来自键盘、文件还是测试数组。这就是把边界留在边缘的价值。

\section{从测试反推接口}

接口不是凭感觉定的。先写出想验证的行为：

\begin{lstlisting}[language=C++,caption={从行为反推接口}]
const std::vector<int> empty;
assert(!summarize_scores(empty).has_value());

const std::vector<int> one{88};
const auto single = summarize_scores(one);
assert(single.has_value());
assert(single->min_score == 88);
assert(single->max_score == 88);
assert(single->average == 88.0);
\end{lstlisting}

这些测试逼着我们承认两个事实：空输入没有统计结果，非空输入必须保留小数平均值。于是 \code{std::optional<ScoreSummary>} 和 \code{double average} 都不是装饰，而是从需求和边界推出来的。

\begin{exercisebox}
把成绩限制为 \code{0} 到 \code{100}。不要在统计函数里偷偷忽略非法值，而是设计一个读取或校验阶段，让非法输入能被明确报告。完成后再给空输入、负数、超过 \code{100}、单元素和普通多元素各写一个测试。
\end{exercisebox}
